\documentclass[submission,copyright,creativecommons]{eptcs}
\providecommand{\event}{LSFA 2026} % Name of the event you are submitting to

\usepackage{iftex}

\ifpdf
  \usepackage{underscore}         % Only needed if you use pdflatex.
  \usepackage[T1]{fontenc}        % Recommended with pdflatex
\else
  \usepackage{breakurl}           % Not needed if you use pdflatex only.
\fi

\title{On the Additive Structure of Quantalic $\lambda$-Calculus}
\author{Renato Neves
\institute{University of Minho and INESC-TEC, Portugal}
\email{nevrenato@di.uminho.pt}
\and
Bruna Salgado
\institute{University of Minho and INESC-TEC, Portugal}
\email{...}
}
\def\titlerunning{On the Additive Structure of Quantalic $\lambda$-Calculus}
\def\authorrunning{Renato Neves and Bruna Salgado}
\begin{document}
\maketitle

\newcommand{\cf}{\emph{cf.}}
\newcommand{\ie}{\emph{i.e.}}
\newcommand{\eg}{\emph{e.g.}}

%% Types
\newcommand{\typefont}[1]{\mathbb{#1}}
\newcommand{\typeOne}{1}
\newcommand{\typeTwo}{2}
\newcommand{\typeA}{\typefont{A}}
\newcommand{\typeB}{\typefont{B}}
\newcommand{\typeC}{\typefont{C}}
\newcommand{\typeV}{\typefont{V}}
\newcommand{\typeD}{\typefont{D}}
\newcommand{\typeE}{\typefont{E}}
\newcommand{\typeF}{\typefont{F}}
\newcommand{\typeI}{\typefont{I}}

%% Sequents
\newcommand{\vljud}{\, \triangleright \,}
\newcommand{\Shuff}{\mathrm{Sf}}
\newcommand{\rulename}[1]{(\textbf{#1})}
\newcommand{\ielim}[2]{#1 \> \> \mathtt{to} \, \ast. \> \> #2}
\newcommand{\telim}[2]{\mathtt{pm} \> \> #1 \> \> \mathtt{to} \> \> x \otimes y. \> \> #2}
\newcommand{\celim}[3]{\mathtt{case} \> \> #1 \> \> \{ \dintrol{x} \Rightarrow #2 ; \>\> 
\dintror{y} \Rightarrow #3 \}  }
\newcommand{\dintrol}[1]{\mathtt{inl}_\typeB (#1)}
\newcommand{\dintror}[1]{\mathtt{inr}_\typeA (#1)}
%% Colimits
\newcommand{\colim}{\mathrm{colim}}
\newcommand{\inl}{\mathrm{inl}}
\newcommand{\inr}{\mathrm{inr}}

%% Semantics
\newcommand{\sem}[1]{\left \llbracket #1 \right \rrbracket}
\newcommand{\prog}[1]{\mathtt{#1}}

\begin{abstract}
Motivated by the lack of methods for reasoning about `case' statements
quantitatively, we introduce \emph{quantalic linear $\lambda$-calculus} with an
additive disjunction operator $\oplus$.  We show that the extended equational
system is sound, and moreover if certain continuity properties (of the
underlying quantale) are adopted that it is (approximately) complete as well. 

We explore different models of the extended calculus, with focus on
probabilistic and quantum programming.  While in the former case we mainly
involve a certain category of Banach spaces, in the latter we recur to the
so-called Karoubi envelope and an enriched reflection between categories of
presheaves.
\end{abstract}

\section{Introduction}
Previous work~\cite{dahlqvist22,dahlqvist2023syntactic} introduced a quantalic
generalisation of linear $\lambda$-calculus, the \emph{exponential-free,
multiplicative fragment} of linear logic. Here we investigate the incorporation
of \emph{additive} structure to this body of work. Our focus is specifically on
the \emph{additive disjunction operator} $\oplus$, which is typically
interpreted via coproducts and gives rise to `case' statements (\ie\
conditionals). Our motivation comes from practice: in trying to reason
quantitatively about (higher-order) programs we often fell short when these
involved conditionals. Of course, as explained below, applications involving
$\oplus$ are broader than this, and fit in the more general pattern of
reasoning quantitatively about co-Cartesian categories enriched over so-called
`\emph{generalised metric spaces}'~\cite{paseka00}.

A number of important results already considered additive structure
in the quantalic setting, even if sometimes implicitly.
References~\cite{mardare2016quantitative,mardare2017axiomatizability,mio24,jurka24}
for example are framed in the setting of universal algebra and therefore
involve additive conjunction (\ie\ $\&$), typically interpreted via categorical
products.  In the higher-order setting, \cite{lago22} enforces additive
conjunction to be left adjoint to implication (interpreted via
Cartesian-closedness), with a series of negative results emerging from this.
Our work is orthogonal to these in that we study the dual of $\&$ (\ie\
$\oplus$) and furthermore we assume the left adjoint of implication to be
multiplicative conjunction (\ie\ $\otimes$) instead of the additive
counterpart. Among other things, this removes the obstacles discussed
in~\cite{lago22}.


Concretely, we extend the quantalic equational system
of~\cite{dahlqvist22,dahlqvist2023syntactic} with the additive disjunction
operator. We show that the extension is sound and complete if certain
\emph{continuity} properties (of the underlying quantale) are adopted.  We also
show that if one drops the well-known \emph{Archimedean
rule}~\cite{mardare2016quantitative,mardare2017axiomatizability} (which
involves infinitely many premisses and is thus problematic) `approximate
completeness' (in the spirit of~\cite{dahlqvist2023syntactic}) is still
retained.

Next, we explore different models of the extended calculus, with focus on
probabilistic~\cite{dahlqvist19,barthe20} and quantum
programming~\cite{selinger2004towards,malherbe13}, two rapidly emerging
computational paradigms. In the former case we will involve the category of
Banach spaces and linear
contractions~\cite{dahlqvist19,dahlqvist2023syntactic}, and illustrate how the
extended framework can be used to reason about Cauchy sequences of
(higher-order) programs -- highlighting the emergent shift from ``seeing
program semantics as the science of program equivalence''~\cite{lago22} to more
flexible, quantitative perspectives, involving functional analysis and beyond.
The quantum case will rely more heavily on categorical machinery: starting with
a certain, enriched quantum category we adopt its \emph{Karoubi
envelope}~\cite{borceux94} to take conditionals into account. Then, in order to
accomodate higher-order structure, we propound a \emph{reflective} subcategory
of the envelope's category of (enriched) presheaves. Remarkably, we also
present a constructive description of the involved reflection by relying on a
\emph{first-order, additive} version of our quantalic calculus (similar in
spirit to the description presented in~\cite{fu22} for the non-enriched
setting).


\bibliographystyle{eptcs}
\bibliography{biblio}
\end{document}
